\section{A zero from one inexact observation per bit}
\label{sec:conditioned-zero}

ProveIt's \texttt{PolynomialBisection} tests the sign of an exact
homogeneous integer numerator. Its companion \texttt{PolynomialRootName}
derives a regular name and a residual schedule for a fixed-degree rational
polynomial. \texttt{PolynomialRootUniqueness} uses a positive lower bound on
a divided difference. For an analytic function a rational argument can have
an irrational value. Waiting for a sign at the midpoint can then fail even at
the first step. Here a lower modulus for increments permits one halving step
with a predetermined precision, including when the midpoint is exactly a zero.

\paragraph{Finite hypotheses.}
Write $e_j=2^{-j}$. At rational $x\in[0,1]$ we have rational observations
$F(x,p)$ with consistency
$F(x,p)-e_p\le F(x,q)+e_q$ for all $p,q$. Define the lower and upper
endpoints $L(x,p)=F(x,p)-e_p$, $U(x,p)=F(x,p)+e_p$.
The statements $f(x)\le0$ and $f(x)\ge0$ abbreviate
$\forall p\,L(x,p)\le0$ and $\forall p\,0\le U(x,p)$, respectively.
Assume $f(0)\le0\le f(1)$. Two numerical witnesses control different parts
of the construction. For all rational sites in the unit interval and all
observation indices the following implications hold:
\begin{align}
 |x-y|\le e_{\mu(k)}
 &\quad\Longrightarrow\quad |F(x,p)-F(y,q)|\le e_k+e_p+e_q,
 \label{eq:zero-forward}\\
 y-x\ge e_k
 &\quad\Longrightarrow\quad
 L(x,p)+e_{g(k)}\le U(y,q).
 \label{eq:zero-inverse}
\end{align}
The second condition is a lower bound on the increase over a specified
distance. A slope bound $f(y)-f(x)\ge2^{-d}(y-x)$ supplies $g(k)=k+d$.
Much flatter functions also qualify. For $f(x)=(x-r)^3$, the identity
\[
 v^3-u^3=(v-u)\left(3((v+u)/2)^2+(v-u)^2/4\right)
\]
supplies $g(k)=3k+2$ for $u=x-r, v=y-r$. This includes a zero with
vanishing derivative. Fixed PR presentations of $F,g,\mu$, possibly with
number parameters, give the first-order reading. The same algorithm runs
relative to observation oracles when those presentations are external.

\begin{theorem}[Conditioned intermediate value]
Under these hypotheses there is a regular dyadic name for the unique
represented zero in $[0,1]$. Starting with $[0,1]$, the first $n$ halvings
make exactly $n$ calls to $F$. The call numbered $j$, beginning at $j=0$,
requests precision $g(j+2)+1$. The forward modulus is used only to certify
the residual. All recursions producing the name are primitive recursive in
the supplied numerical data.
\end{theorem}

\begin{proof}
At stage $j$ write $I=[a,b]$, $w=b-a=e_j$, and
$m=(a+b)/2$, $u=m-w/4$, $v=m+w/4$. Set $p=g(j+2)+1$, $h=e_p$ and obtain
$z=F(m,p)$. Form the next bracket by the rational tests
\[
 [a,m]\quad(z\ge h),\qquad [m,b]\quad(z\le-h),\qquad
 [u,v]\quad(-h<z<h).
\]
In the first case consistency gives $0\le z-h\le U(m,t)$ at every
test precision $t$. In the second case $L(m,t)\le z+h\le0$.
In the remaining case both gaps $m-u=v-m=e_{j+2}$ qualify for
\eqref{eq:zero-inverse}. Since $e_{g(j+2)}=2h$,
\[
 L(u,t)+2h\le z+h,\qquad z-h+2h\le U(v,t).
\]
These two inequalities certify the signs of the new endpoints. Every
branch gives a nested bracket of width $w/2$. For the induction in iPRA,
fix an arbitrary $t$ first. At this $t$, the two sign inequalities for the
current endpoints are a quantifier-free invariant; the transition uses the
gap contract instantiated at $(p,t)$ or $(t,p)$. Generalize $t$ only after
induction. The separate nesting and width invariants are rational
inequalities or identities.

Let $m_j$ be the midpoint of the $j$th bracket. For $d\ge0$,
$m_{j+d}$ belongs to that bracket, so $|m_{j+d}-m_j|\le e_j/2$.
This gives a regular name. Fix $k$, take $t=k+3$, and take any
$j\ge\mu(k+1)$. Applying \eqref{eq:zero-forward} between the midpoint and
each endpoint yields
\[
 |F(m_j,t)|\le e_{k+1}+3e_t=7e_k/8.
\]
Indeed the left endpoint has $F(a,t)\le e_t$ and the right endpoint has
$F(b,t)\ge-e_t$. The former bounds the midpoint observation from above,
the latter from below. These tests at arbitrarily prescribed precision say
that the completed evaluation is zero. Explicitly, if an evaluation
convention uses $m_i$ and an observation at $q$, choose a later $j\ge i$.
Whenever $e_i/2\le e_{\mu(s)}$, forward continuity bounds
$|F(m_i,q)-F(m_j,t)|\le e_s+e_q+e_t$. Together with the displayed
residual, this supplies the comparison to zero with any desired slack.
All stage choices are available from the numerical moduli.

For uniqueness, suppose two represented zeros have points at positive
distance. In the finite-test formulation it suffices to rule out a rational
apartness witness in either order. Such a witness supplies $k$ and a
precision beyond which the two streams of rational sites, in the indicated
order, are at least $e_k$ apart (shrink the original gap by a factor two).
Evaluate both streams sufficiently far out that their absolute observations
are at most $e_{g(k)}/8$, with observation radii at most that same amount.
Their upper-to-lower difference is at most $e_{g(k)}/2$, contradicting
\eqref{eq:zero-inverse}. For the equality convention based on closed
rational inequalities, a failed finite equality test itself gives the
rational apartness witness, by consistency. Decidability of that one test
therefore gives equality at every precision.

Extensionality also has a direct finite proof. Brackets from two valid
observations of the same $f$ must intersect. If a positive endpoint of one
were strictly to the left of a negative endpoint of the other, choose $k$
with $e_k$ below their positive rational distance and then $p$ with
$4e_p<e_{g(k)}$. Transfer the endpoint signs to one presentation using
consistency across the two names. Specifically, a common represented value
with a nonnegative sign satisfies $F'(x,p)\ge-e_p$ in either presentation:
the cross bound $F(x,t)\le F'(x,p)+e_t+e_p$ and $F(x,t)\ge-e_t$
give $F'(x,p)\ge-e_p-2e_t$ for every $t$; any failure at $p$ is
refuted by choosing $2e_t$ below the rational gap. The nonpositive sign transfers
by the reversed calculation. Now \eqref{eq:zero-inverse} gives
$e_{g(k)}\le4e_p$, a contradiction. Rational endpoints are decidably
ordered, so intersection follows. Midpoints of intersecting intervals
differ by at most the sum of their half-widths. Thus alternate rounding or
choices of valid witnesses give equivalent names even when their branch
sequences differ.
\end{proof}

\paragraph{An example with genuinely inexact midpoint values.}
Take $f(x)=\exp(x)-3/2$ on the unit interval and define its reading at
$(x,p)$ by $\sum_{i=0}^{p+2}x^i/i!-3/2$. For $i\ge1$ one has
$i!\ge2^{i-1}$; hence the tail after $M=p+2$ is at most
$\sum_{i>M}2^{1-i}=2^{1-M}<e_p$. For $y\ge x$ all positive-index
monomial increments are nonnegative and the linear term contributes
$y-x$, giving $g(k)=k$. The finite identity for $y^i-x^i$ bounds each
increment by $i(y-x)$. The resulting total Lipschitz bound is at most
$\sum_{i\ge1}i/i!=\sum_{i\ge0}1/i!\le3$, so $\mu(k)=k+2$ suffices.
The first two terms at $1$ and the exact value at $0$ give the required
signs. The algorithm names $\log(3/2)$ by evaluating rational Taylor
polynomials with stage-$j$ observation index $j+3$. It needs neither a
logarithm routine nor an exact sign test for an exponential.

\paragraph{Storage and verification.}
Write the left endpoint in stage $j$ as $A/2^{j+1}$, so the right numerator
at that scale is $A+2$. The next left numerator, at denominator $2^{j+2}$,
is $2A$, $2A+2$, or $2A+1$ for the left, right, or central branch.
One integer of $O(j+1)$ bits and the stage number suffice between calls.
The total running time also includes the cost and output lengths of the
evaluations; a count of calls alone gives no general polynomial-time claim.
\texttt{Zeros.lean} implements the three branches on \texttt{Rat}.
\texttt{goodAt\_unit} derives the fixed-precision bracket invariant from
the global separation hypothesis. \texttt{width\_run} and
\texttt{later\_midpoints} prove the width and Cauchy estimates;
\texttt{residual\_unit} derives the full scheduled $7e_k/8$ residual bound
for every later stage from the forward and inverse moduli.
\texttt{no\_cross} checks the scalar separation contradiction.

\texttt{ZeroObject.lean} also proves a closed theorem in the repository's
ordinary object calculus for the three-way guard. Multiply the finitely
many rational observations and their radii by a common positive denominator.
Translate the scaled signed observations by a common integer offset $K$
large enough to make them nonnegative. Write $C$ for the translated midpoint
observation, $h$ for its scaled radius, $L_m,U_m$ for the translated bounds
at midpoint test precision $t$, and $L_u,U_v$ for the translated quarter-point
bounds. The instantiated consistency and gap premises become
\[
 L_m\le C+h,\quad C\le U_m+h,\quad L_u+h\le C,\quad C+h\le U_v.
\]
The object theorem tests $K+h\le C$ first and $C+h\le K$ second.
It pairs each of the three disjoint guards with the required signs
$K\le U_m$, $L_m\le K$, or $L_u\le K\le U_v$, respectively.
For example the third case cancels $h$ from
$L_u+h\le C<K+h$. The seven variables in this sequent are universally
closed in \texttt{branches\_closed}. To extend this local object theorem
to the full stage recursion, use the integer
state $A$ above, encode pairs of signed observations by positive
denominators, and clear those denominators in each occurrence of $L,U$.
Implement both rational comparisons by reflected weak integer order.
Instantiate \eqref{eq:zero-forward} and \eqref{eq:zero-inverse} at the
indicated rational sites and precisions. Apply the repository's
quantifier-free trace induction to the conjunction of the two sign tests
for fixed $t$, the numerator range, and the denominator identity.
For the width over a window fix the starting stage as another free
parameter. The remaining equality tests for completed names follow by the
displayed rational-gap arguments. The included Lean module checks the
rational model; this last translation of the full case study into
\texttt{PTm} derivations is a formalization plan.

\paragraph{A residual witness without the inverse modulus.}
The same branches yield an approximate intermediate value theorem for any
function satisfying \eqref{eq:zero-forward} and the endpoint signs.
Given $k$, hold $p=k+3$ fixed and perform at most $\mu(k+1)$ midpoint queries.
For a decisive sign take the corresponding half. At the first ambiguous
sample return that midpoint $x$. Consistency gives
$|F(x,t)|\le 2e_p+e_t$ at every $t$. If all tests are decisive, return
the midpoint of the final bracket of width $e_{\mu(k+1)}$.
Using its endpoint signs at precision $p$ in \eqref{eq:zero-forward}
gives $|F(x,t)|\le e_{k+1}+2e_p+e_t<e_k+e_t$ for every $t$.
The bound covers $\mu(k+1)=0$, when the initial midpoint is returned.
This yields a PR rational witness with certified residual at each
prescribed accuracy. Independent calls need not produce a Cauchy name.
It also replaces an exhaustive grid of up to $2^{\mu(k+1)}+1$ points
by at most $\mu(k+1)$ scalar queries.

\paragraph{Two obstructions at the boundary.}
There is no uniform query-precision ceiling, even for an error target of
$1/8$, for all strictly increasing functions given only by observations,
the forward Lipschitz bound $1$, and the weak endpoint signs. For any
proposed ceiling $P$, take $f_r(x)=e_{P+1}(x-r)$ for
$r=1/4$ or $3/4$. The reply $0$ is valid for both at every site and every
precision up to $P$. An algorithm receiving these common replies must
produce the same rational approximation, yet the permitted approximations
for the two roots are disjoint. Both functions even have strict endpoint signs.
This argument rules out a uniform ceiling on this representation; on strictly
increasing inputs an adaptive search can request further precision.

Allowing a zero interval creates a stronger obstruction. Put
\[
 f_0(x)=\min(x-1/4,0)+\max(x-3/4,0).
\]
It is nondecreasing and $1$-Lipschitz. For rational $0<h<1/8$,
$f_0+h$ has the unique root $1/4-h$ and $f_0-h$ has the unique root
$3/4+h$. The endpoint signs for both have margin at least $1/8$.
Run any proposed total exact zero-name realizer on the exact rational
evaluator for $f_0$ until it produces an answer of radius $1/8$.
This finite transcript uses some greatest precision $P$ (use $P=0$ for
an empty transcript). Choose $h<\min(1/8,e_P)$.
The same answers at those queries occur in valid names for both perturbed
functions: use the original answers for those finitely many pairs and
exact answers elsewhere. Continuity and endpoint promises still hold.
The two roots differ by more than $1/2$, so one answer of radius $1/8$
cannot serve both. This is a finite-observation impossibility proof for a
total extensional root selector on this representation. It leaves the
usual approximate intermediate value theorem intact.
